\chapter{ВЫБОР И ОБОСНОВАНИЕ НАПРАВЛЕНИЯ ПРОЕКТИРОВАНИЯ И ИСПОЛЬЗУЕМЫХ МЕТОДИК}

\section{Методология проектирования программно-аппаратного комплекса}

Проектирование системы мониторинга базируется на итеративном подходе, включающем этапы моделирования, макетирования и верификации. Согласно п. 1.2.12.2 СТП 01-2022, методология сбора данных основывается на событийной модели передачи сообщений (event-driven architecture).

В качестве основного метода анализа достоверности передаваемой информации используется статистическая обработка выборок с применением фильтра скользящего среднего для отсечения кратковременных помех датчиков.

\section{Обоснование выбора аппаратной платформы}

Для реализации узла сбора данных рассматривались микроконтроллеры архитектур AVR, ARM и Xtensa. Сравнительный анализ по критериям энергопотребления в режиме «Deep Sleep» и наличию встроенных стеков беспроводной связи приведен на рисунке 2.1.

% Пример вставки рисунка (файл image.jpg из вашего списка)
\begin{figure}[h!]
    \centering
    \includegraphics[width=0.7\textwidth]{image.jpg}
    \caption{Структурная схема взаимодействия аппаратных модулей системы}
    \label{fig:structure}
\end{figure}

На основании анализа выбран микроконтроллер ESP32 (архитектура Xtensa LX6), обладающий следующими преимуществами:
\begin{itemize}
    \item встроенный модуль Wi-Fi/Bluetooth;
    \item наличие двух процессорных ядер;
    \item поддержка аппаратного шифрования данных AES-256.
\end{itemize}

\section{Выбор программного стека и протоколов обмена}

Для обеспечения отказоустойчивости системы на уровне программного обеспечения выбрана операционная система реального времени FreeRTOS. Это позволяет эффективно распределять задачи между ядрами микроконтроллера:
\begin{enumerate}
    \item ядро 0 --- обслуживание сетевого стека и поддержание Wi-Fi соединения;
    \item ядро 1 --- опрос датчиков и первичная обработка сигналов.
\end{enumerate}

В качестве протокола прикладного уровня выбран MQTT версии 3.1.1 с использованием механизма «Last Will and Testament» для оперативного уведомления сервера о потере связи с узлом мониторинга.

\section{Математическая модель оценки энергопотребления}

Средний ток потребления устройства $I_{avg}$ рассчитывается по формуле \eqref{eq:current}:

\begin{equation}
    I_{avg} = \frac{I_{active} \cdot t_{active} + I_{sleep} \cdot t_{sleep}}{t_{active} + t_{sleep}}
    \label{eq:current}
\end{equation}

\begin{formula_where}
    $I_{active}$ --- ток в активном режиме передачи данных, мА; \\
    $t_{active}$ --- время нахождения в активном режиме, с; \\
    $I_{sleep}$ --- ток в режиме глубокого сна, мА; \\
    $t_{sleep}$ --- период нахождения в режиме сна, с.
\end{formula_where}

Расчет по данной модели позволяет прогнозировать срок службы устройства от одного литиевого аккумулятора емкостью 2500 мАч.

\newpage